% =========================================================================
\section{Funciones de Green}\label{sec:green}
% =========================================================================

% -------------------------------------------------------------------------
\subsection{Idea: descomponer la fuente}\label{sec:idea-green}
% -------------------------------------------------------------------------
Por el principio de superposición (\cref{prop:linealidad}), si
$b=\sum_kb_k$ y $\Lop[y_{p,k}]=b_k$, entonces $y_p=\sum_ky_{p,k}$: el efecto
de una combinación lineal de fuentes es la superposición del efecto de cada
una. ¿Cuál es la forma más elemental de trocear una fuente? La respuesta
son las \emph{fuentes puntuales}. Por la propiedad de filtrado de la delta
de Dirac,
\begin{equation}\label{eq:filtrado}
  b(x)=\int b(\tilde x)\,\delta(x-\tilde x)\,\dif\tilde x .
\end{equation}
La fuente se ve así como una fila de picos infinitamente estrechos: en cada
punto $\tilde x$ hay una fuente puntual $\delta(x-\tilde x)$ de intensidad
$b(\tilde x)\,\dif\tilde x$, y la integral es la versión continua de la suma
$\sum_kb_k$ (el índice discreto $k$ pasa a ser la variable continua
$\tilde x$). Basta, pues, resolver el problema para un impulso unitario
(\cref{fig:superposicion}).

% Figura: fuente troceada en impulsos y superposición de respuestas
% (datos: scripts/generar_datos.py; operador y''+2*0.16 y'+y)
\begin{figure}[ht]
\centering
\begin{tikzpicture}
  \begin{groupplot}[group style={group size=1 by 2, vertical sep=1.45cm},
      oscuro, width=0.94\textwidth, xmin=0, xmax=20, xtick={0,2,...,20}]
    % ---------- (a) fuente e impulsos ----------
    \nextgroupplot[height=4.0cm, ymin=-0.05, ymax=1.2, ytick={0,0.5,1},
        title={(a) Fuente $b(\tilde x)$ troceada en $K=8$ impulsos de anchura $\Delta\tilde x=0{,}7$}]
      \addplot[blanco, line width=1.3pt] table[x=x, y=b] {datos/superposicion.dat};
      \foreach \k/\s/\c in {1/0.35/cm1,2/1.05/cm2,3/1.75/cm3,4/2.45/cm4,
                            5/3.15/cm5,6/3.85/cm6,7/4.55/cm7,8/5.25/cm8}{
        \edef\temp{\noexpand\draw[-{Stealth[length=4pt,width=4pt]}, line width=1.3pt, \c]
          (axis cs:\s,0) -- (axis cs:\s,{sin(deg(pi*\s/5.6))^2});}
        \temp}
      \node[font=\scriptsize, text=gris, anchor=west] at (axis cs:6.4,0.75)
        {impulsos $b(\tilde x_k)\,\Delta\tilde x\;\delta(x-\tilde x_k)$};
    % ---------- (b) respuestas y superposición ----------
    \nextgroupplot[height=5.4cm, ymin=-0.85, ymax=1.5, ytick={-0.5,0,0.5,1},
        xlabel={$x$},
        title={(b) Respuestas $b(\tilde x_k)\,G(x,\tilde x_k)\,\Delta\tilde x$ y su suma},
        legend style={at={(0.985,0.96)}, anchor=north east}]
      \addplot[line width=0.8pt, cm1, forget plot] table[x=x, y=r1] {datos/superposicion.dat};
      \addplot[line width=0.8pt, cm2, forget plot] table[x=x, y=r2] {datos/superposicion.dat};
      \addplot[line width=0.8pt, cm3, forget plot] table[x=x, y=r3] {datos/superposicion.dat};
      \addplot[line width=0.8pt, cm4, forget plot] table[x=x, y=r4] {datos/superposicion.dat};
      \addplot[line width=0.8pt, cm5, forget plot] table[x=x, y=r5] {datos/superposicion.dat};
      \addplot[line width=0.8pt, cm6, forget plot] table[x=x, y=r6] {datos/superposicion.dat};
      \addplot[line width=0.8pt, cm7, forget plot] table[x=x, y=r7] {datos/superposicion.dat};
      \addplot[line width=0.8pt, cm8, forget plot] table[x=x, y=r8] {datos/superposicion.dat};
      \addplot[ambar, line width=1.2pt, densely dashed] table[x=x, y=suma] {datos/superposicion.dat};
      \addplot[blanco, line width=1.6pt] table[x=x, y=yp] {datos/superposicion.dat};
      \legend{suma de las $8$ respuestas, $y_p=\int G\,b\,\dif\tilde x$}
  \end{groupplot}
\end{tikzpicture}
\caption{La función de Green como respuesta a impulsos, para
$y''+2\gamma y'+y=b$ con $\gamma=0{,}16$. (a) La fuente
$b(\tilde x)=\sin^{2}(\pi\tilde x/5{,}6)$ se aproxima por $8$ impulsos
(flechas). (b) Cada impulso produce una respuesta proporcional a
$G(x,\tilde x_k)$ \eqref{eq:green-amortiguado}, nula antes de $\tilde x_k$
(causalidad); su suma (trazo discontinuo) aproxima la solución
$y_p=\int G\,b\,\dif\tilde x$ (trazo continuo), a la que tiende cuando
$\Delta\tilde x\to0$.}
\label{fig:superposicion}
\end{figure}


% -------------------------------------------------------------------------
\subsection{Definición y fórmula de la solución}\label{sec:def-green}
% -------------------------------------------------------------------------

\begin{definicion}[Función de Green]\label{def:green}
Una \emph{función de Green} del operador $\Lop$ es una función
$G(x,\tilde x)$, $x,\tilde x\in I$, que cumple
\begin{equation}\label{eq:green}
  \Lop_x\,G(x,\tilde x)=\delta(x-\tilde x),
\end{equation}
donde el subíndice indica que $\Lop$ deriva respecto de $x$ (el punto donde
se \emph{observa} el efecto) y $\tilde x$ (el punto donde está la
\emph{fuente}) actúa como un parámetro. La igualdad se entiende en sentido
distribucional: $\Lop_xG=0$ para $x\neq\tilde x$, con las condiciones de
salto en $x=\tilde x$ de la \cref{prop:salto}
\textup{\cite{STAKGOLD}, \cite[cap.~10]{ARFKEN}}.
\end{definicion}

\begin{teorema}[Solución como superposición de impulsos]\label{teo:representacion}
Si $G$ es una función de Green de $\Lop$, entonces
\begin{equation}\label{eq:representacion}
  \boxed{\;y_p(x)=\int G(x,\tilde x)\,b(\tilde x)\,\dif\tilde x\;}
\end{equation}
es una solución particular de $\Lop[y]=b$.
\end{teorema}

\begin{proof}
A la fuente $b(\tilde x)\,\delta(x-\tilde x)$ le corresponde la respuesta
$b(\tilde x)\,G(x,\tilde x)$; sumando sobre $\tilde x$ y como $\Lop_x$ sólo
actúa sobre $x$, entra en la integral y sólo afecta a $G$:
\begin{equation}\label{eq:comprobacion}
  \Lop_x\,y_p(x)=\int\ubr{\Lop_xG(x,\tilde x)}{\delta(x-\tilde x)}\,b(\tilde x)\,\dif\tilde x=b(x).
\end{equation}
Para la función de Green causal, el \cref{teo:green-variacion} reduce este
cálculo formal al \cref{teo:variacion}, ya demostrado.
\end{proof}

\begin{observacion}[$G$ depende del operador, no de la fuente]\label{obs:operador}
$G$ está asociada a $\Lop$ (y a las condiciones de contorno o iniciales que
se impongan), no a $b$: se calcula una sola vez y sirve para cualquier
fuente. Físicamente, $G(x,\tilde x)$ es la \emph{respuesta en $x$ a un
impulso unitario en $\tilde x$}, y la solución es la suma ponderada de todas
esas respuestas.
\end{observacion}

% -------------------------------------------------------------------------
\subsection{No unicidad, causalidad y condiciones de salto}\label{sec:salto}
% -------------------------------------------------------------------------
$G$ no es única: si $h(x;\tilde x)$ cumple $\Lop_xh=0$, también $G+h$
satisface \eqref{eq:green}. Se fija imponiendo condiciones de contorno o
iniciales. Una elección habitual es la \emph{causal} o \emph{retardada}:
\begin{equation}\label{eq:causal}
  G(x,\tilde x)=0\qquad\text{para } x<\tilde x
\end{equation}
(no hay respuesta antes del impulso).

\begin{proposicion}[Condiciones de salto]\label{prop:salto}
Sea $G$ de clase $\mathcal C^{n}$ en $x\neq\tilde x$ y solución de
\eqref{eq:green}. Entonces $G,\partial_xG,\dots,\partial_x^{\,n-2}G$ son
continuas en $x=\tilde x$ y
\begin{equation}\label{eq:salto}
  \partial_x^{\,n-1}G\,\Big|_{x=\tilde x^-}^{x=\tilde x^+}=\frac{1}{a_n(\tilde x)} .
\end{equation}
\end{proposicion}

\begin{proof}[Justificación]
La delta sólo puede equilibrarse con la derivada de mayor orden: si
$\partial_x^{\,n-1}G$ salta en $c$, entonces $\partial_x^{\,n}G$ contiene
$c\,\delta(x-\tilde x)$, y $a_n\partial_x^{\,n}G$ aporta
$a_n(\tilde x)\,c\,\delta$. Si alguna derivada de orden $j\le n-2$ saltara,
$\partial_x^{\,n}G$ contendría $\delta^{(n-1-j)}$, que nada en $\Lop_xG$
podría compensar. Integrando \eqref{eq:green} en
$(\tilde x-\varepsilon,\tilde x+\varepsilon)$ y haciendo $\varepsilon\to0$,
sólo sobrevive $a_n(\tilde x)\,c=1$.
\end{proof}

% Figura: función de Green causal para x~ = 2 (continuidad y salto de la derivada)
\begin{figure}[ht]
\centering
\begin{tikzpicture}
  \begin{axis}[oscuro, width=0.94\textwidth, height=5.8cm,
      xmin=0, xmax=16, ymin=-1.15, ymax=1.62, xtick={0,2,...,16},
      ytick={-1,-0.5,0,0.5,1}, xlabel={$x$},
      samples=500, domain=0:16, legend pos=north east]
    \fill[grisoscuro, opacity=0.45] (axis cs:0,-1.15) rectangle (axis cs:2,1.62);
    \addplot[c1] {x<2 ? 0 : sin(deg(x-2))};
    \addplot[c2] {x<2 ? 0 : exp(-0.16*(x-2))*sin(deg(0.987117*(x-2)))/0.987117};
    % pendientes a ambos lados de x~ = 2
    \draw[ambar, line width=0.9pt, densely dashed] (axis cs:2,0) -- (axis cs:3.05,1.05);
    \draw[ambar, line width=0.9pt, densely dashed] (axis cs:0.95,0) -- (axis cs:2,0);
    \addplot[only marks, mark=*, mark size=2pt, ambar, forget plot] coordinates {(2,0)};
    \legend{$\gamma=0$: $\Heav(x-\tilde x)\sin(x-\tilde x)$,
            $\gamma=0{,}16$: amortiguada}
    \node[font=\scriptsize, text=gris, align=center, anchor=south] at (axis cs:1,-1.1)
      {$G=0$\\(causalidad)};
    \node[font=\scriptsize, text=ambar, anchor=south] at (axis cs:2.55,1.25)
      {pendiente $1$};
    \node[font=\scriptsize, text=ambar, anchor=south] at (axis cs:1.05,0.06)
      {pendiente $0$};
    \node[font=\small, text=texto, anchor=north] at (axis cs:2,-0.08) {$\tilde x$};
  \end{axis}
\end{tikzpicture}
\caption{Función de Green causal de $y''+2\gamma y'+y=\delta(x-\tilde x)$ con
$\tilde x=2$, sin amortiguamiento \eqref{eq:green-oscilador} y con él
\eqref{eq:green-amortiguado}. $G$ es nula antes del impulso (zona
sombreada) y continua en $x=\tilde x$, pero su derivada salta de $0$ a
$1=1/a_2$, como exige la \cref{prop:salto}: el impulso comunica
«velocidad» instantáneamente, no «posición».}
\label{fig:green}
\end{figure}


% -------------------------------------------------------------------------
\subsection{Conexión con la variación de las constantes}\label{sec:conexion}
% -------------------------------------------------------------------------

\begin{teorema}[Función de Green causal]\label{teo:green-variacion}
La función de Green causal de $\Lop$ es
\begin{equation}\label{eq:green-causal}
  \boxed{\;G(x,\tilde x)=\Heav(x-\tilde x)\;\yv^{\T}(x)\,W^{-1}(\tilde x)\,\en\,\frac{1}{a_n(\tilde x)}\;}
\end{equation}
Es decir: para $x>\tilde x$, $G(\cdot,\tilde x)$ es la solución de la
homogénea con datos $\bigl(0,\dots,0,1/a_n(\tilde x)\bigr)$ en $\tilde x$.
Además, \eqref{eq:representacion} con $G$ causal y fuente en
$[x_0,\infty)$ coincide con \eqref{eq:variacion}.
\end{teorema}

\begin{proof}
Para $x>\tilde x$, \eqref{eq:green-causal} es combinación lineal de las
$\hat y_i$, luego $\Lop_xG=0$; para $x<\tilde x$ es nula. En
$x=\tilde x^{+}$,
\[
  \partial_x^{\,j}G\big|_{x=\tilde x^+}
  =\bigl(\yv^{(j)}(\tilde x)\bigr)^{\T}W^{-1}(\tilde x)\,\en\,\frac{1}{a_n(\tilde x)}
  =\frac{\delta_{j,n-1}}{a_n(\tilde x)},\qquad j=0,\dots,n-1,
\]
porque $\bigl(\yv^{(j)}\bigr)^{\T}$ es la fila $j$ de $W$ y $W W^{-1}=\mathbb 1$.
Así se cumplen las condiciones de la \cref{prop:salto}. Finalmente,
introduciendo $\yv^{\T}(x)$, que no depende de $s$, dentro de la integral
de \eqref{eq:variacion},
\begin{equation}\label{eq:variacion-green}
  y_p(x)=\int_{x_0}^{x}\ubr{\yv^{\T}(x)\,W^{-1}(s)\,\en\,\frac{1}{a_n(s)}}{G(x,s)}\;b(s)\,\dif s ,
\end{equation}
y el límite superior $x$ expresa la causalidad: sólo contribuyen las fuentes
con $s<x$.
\end{proof}

La variación de las constantes \emph{es} el método de Green: son el mismo
cálculo visto desde dos lados.

\begin{corolario}[Orden $2$]\label{cor:green-n2}
Para $n=2$,
\begin{equation}\label{eq:green-n2}
  G(x,\tilde x)=\Heav(x-\tilde x)\,
  \frac{\hat y_1(\tilde x)\,\hat y_2(x)-\hat y_1(x)\,\hat y_2(\tilde x)}{a_2(\tilde x)\,w(\tilde x)} .
\end{equation}
\end{corolario}

\begin{ejemplo}[Oscilador]\label{ej:green-oscilador}
Para $y''+\omega^{2}y=\delta(x-\tilde x)$, con $\hat y_1=\cos\omega x$,
$\hat y_2=\sin\omega x$, $w=\omega$ y $a_2=1$,
\begin{equation}\label{eq:green-oscilador}
  G(x,\tilde x)=\Heav(x-\tilde x)\,\frac{\sin\omega(x-\tilde x)}{\omega},
\end{equation}
que sustituida en \eqref{eq:representacion} da exactamente
\eqref{eq:forzado}. $G$ es la respuesta de un oscilador en reposo a un golpe
unitario en el instante $\tilde x$: a partir de ahí oscila como
$\sin\omega(x-\tilde x)$.
\end{ejemplo}

\begin{ejemplo}[Oscilador amortiguado]\label{ej:green-amortiguado}
Para $y''+2\gamma y'+\omega^{2}y=\delta(x-\tilde x)$ con $0<\gamma<\omega$
($\gamma=\nu\omega/2$ en la notación del \cref{ej:amortiguado}), la base
$\hat y_{1,2}=\eul^{-\gamma x}\{\cos\Omega x,\sin\Omega x\}$,
$\Omega=\sqrt{\omega^{2}-\gamma^{2}}$, tiene $w=\Omega\,\eul^{-2\gamma x}$, y
\eqref{eq:green-n2} da
\begin{equation}\label{eq:green-amortiguado}
  G(x,\tilde x)=\Heav(x-\tilde x)\,\eul^{-\gamma(x-\tilde x)}\,\frac{\sin\Omega(x-\tilde x)}{\Omega} .
\end{equation}
\end{ejemplo}

\begin{observacion}[Invariancia por traslación]\label{obs:convolucion}
Con coeficientes constantes, $G$ depende sólo de $x-\tilde x$:
$G(x,\tilde x)=g(x-\tilde x)$. Entonces \eqref{eq:representacion} es una
\emph{convolución}, $y_p=g*b$.
\end{observacion}

% -------------------------------------------------------------------------
\subsection{Puente con el Electromagnetismo}\label{sec:poisson}
% -------------------------------------------------------------------------
El mismo esquema funciona en tres dimensiones. En electrostática el
operador es el laplaciano y la fuente, la densidad de carga: la ecuación de
Poisson es $\nabla^{2}\phi=-\rho/\varepsilon_0$. La función de Green del
laplaciano en todo el espacio, con $G\to0$ en el infinito, es
\begin{equation}\label{eq:green-laplaciano}
  \nabla_{\bm r}^{2}\,G(\bm r,\bm r')=-\delta^{3}(\bm r-\bm r'),
  \qquad
  G(\bm r,\bm r')=\frac{1}{4\pi\abs{\bm r-\bm r'}},
\end{equation}
y \eqref{eq:representacion}, en tres dimensiones, da
\begin{equation}\label{eq:coulomb}
  \phi(\bm r)=\frac1{\varepsilon_0}\int G(\bm r,\bm r')\,\rho(\bm r')\,\dif^{3}r'
            =\frac{1}{4\pi\varepsilon_0}\int\frac{\rho(\bm r')}{\abs{\bm r-\bm r'}}\,\dif^{3}r' .
\end{equation}
Físicamente, $G$ es el potencial de una carga puntual unidad en $\bm r'$, y
\eqref{eq:coulomb} es la ley de Coulomb: el potencial de una distribución
es la superposición de los potenciales de las cargas puntuales
$\dif q=\rho\,\dif^{3}r'$ \cite[§2.3]{GRIFFITHS}. (El signo de
\eqref{eq:green-laplaciano} es una convención habitual en electrostática.)
